\section{Transformações}

\begin{def:transformacao}
Uma transformação de uma variável aleatória $X: \Omega \to \mathbb{R}$ com imagem $\mathcal{X}$
é toda função $f: \mathcal{X} \to \mathbb{R}$.
\end{def:transformacao}

Toda transformação define uma nova variável aleatória.
Isso porque a função composta $f \circ X: \Omega \to \mathbb{R}$ é, por definição,
uma variável alaeatória.

\begin{def:variavel derivada}
Se $X: \Omega \to \mathbb{R}$ é uma variável aleatória e $f$ é uma transformação de $X$,
então $f(X)$ equivale à variável definida por
\[
	f \circ X
\]
\end{def:variavel derivada}

A transformação não necessariamente é bijetora, o que nos leva a ter de tratar sobre a
sua distribuição de probabilidade a partir da distribuição de $X$ com um pouco mais de
cuidado. Antes de tudo, iremos definir uma generalização de uma função inversa para
$f$.

\begin{def:generalizacao inversa}
Se $f: D \to \mathbb{R}$ é uma função real, então $f^\dagger: \Parts{\mathbb{R}} \to \Parts{D}$ é
tal que
\[
	f^\dagger(A) = \{x \in D \mid f(x) \in A\}.
\]
Em especial, se $y \in \mathbb{R}$, então
\[
	f^\dagger(y) = f^\dagger(\{y\})
\]
\end{def:generalizacao inversa}

\begin{prop:prob transformacao}
\label{prop:prob transformacao}
Seja $X$ uma variável aleatória, $f$ uma transformação e $A \subset \mathbb{R}$.
\[
	P(f(X) \in A) = P(X \in f^\dagger(A))
\]
\begin{proof}
	\begin{align*}
		P(f(X) \in A) & = P(\{x \in \mathbb{R} \mid f(x) \in A\}) \\
		              & = P(X \in f^\dagger(A))
	\end{align*}
\end{proof}
\end{prop:prob transformacao}

\subsection{Transformação de Variáveis Discretas}
\begin{thm:prob transformacao discreta}
Seja $X$ uma variável discreta com suporte $\Omega_X$, f.m.p. igual a $p_X$,
$f$ uma transformação de $X$ e $Y = f(X)$.
A variável $Y$ é discreta e sua f.m.p é dada por
\[
	p_Y(y) = \sum_{x \in f^\dagger(y)} p_X(x)
\]
\begin{proof}
	De \ref{prop:prob transformacao} sabemos que
	\[
		P(Y \in A) > 0 \Longleftrightarrow P(X \in f^\dagger(A)) > 0,
	\]
	o que ocorre também se, e somente se, $A \cap \Omega_X \neq \varnothing$.
	Essa afirmação é verdade somente quando há $x \in \mathbb{R}$ que satisfaz $f(x) \in A$
	e $x \in \Omega_X$.
	Portanto, $y \in \Omega_Y$ ocorre somente quando $y = f(x)$ e $x \in \Omega_X$, ou
	seja, $\Omega_Y$ é a imagem de $\Omega_X$ sobre $f$, sendo, pois, enumerável.

	O resultado sobre a f.m.p. de $Y$ segue do fato de que
	\[
		p_Y(y) = P(Y = y) = P(X \in f^\dagger(y))
	\]
	e pela proposição \ref{prop:fmp}, temos que
	\[
		p_Y(y) = \sum_{x \in f^\dagger(y)} p_x(X)
	\]
\end{proof}
\end{thm:prob transformacao discreta}

\subsection{Transformação de Variáveis Contínuas}

Para transformações de variáveis contínuas, concentramos-nos, primeiramente, nas
transformações mńótonas, pois essas são bijetoras, e compõem casos mais simples para
se analisar.

É possível obter a FDA de uma transformação monótona $Y = g(X)$ de uma variável $X$ apenas
em termos da própria FDA de $X$ e da inversa da transformação.

\begin{thm:fda transformacao continua}
Seja $X$ uma variável aleatória com FDA igual a $F_X$ e suporte $\Omega_X$.
Seja também $Y=g(X)$ uma transformação monótona com suporte $\Omega_Y$.
Se $g$ for crescente, então
\[
	\forall x \in \Omega_{X} \forall y \in \Omega_{Y},\ F_{Y}(y) =  F_{X}(g^{-1}(y)).
\]
Se for decrescente, então
\[
	\forall x \in \Omega_{X} \forall y \in \Omega_{Y},\ F_{Y}(y) = 1 - F_{X}(g^{-1}(y)).
\]
\end{thm:fda transformacao continua}

Por sua vez, a densidade de transformações $Y=g(X)$ em casos em que a densidade de $X$
é contínua, pode ser obtidas via a Regra da Cadeia do Cálculo Diferencial, com a
condição de que a inversa de $g$ seja também contínua em $\Omega_Y$.
Caso a densidade de $X$ nem a inversa da transoformação sejam contínuas em $\Omega_X$ e
$\Omega_Y$ respectivamente, a regra da cadeia não poderia ser aplicada, complexificando
a determinação da densidade de $Y$.

\begin{thm:fdp transformacao continua}
\label{thm:fdp transformacao continua}
Seja $X$ uma variável aleatória com f.d.p. igual a $f_X$.
Seja $Y=g(X)$ uma variável obtida pela transformação $g$ e com suporte $\Omega_Y$.
Se $g$ é monótona, $f_X$ é contínua em $\Omega_X$ e $g^{-1}$ tem derivada contínua em
$\Omega_Y$, então a f.d.p. de $Y$ é dada por
\[
	f_{Y}(y) = \begin{cases}
		f_{X}(g^{-1}(y)) \cdot \left| \frac{d}{dy} g^{-1}(y)\right|, & \text{ se } y \in \Omega_{Y} \\
		0, \text{ caso contrário.}
	\end{cases}
\]
\end{thm:fdp transformacao continua}

O teorema \ref{thm:fdp transformacao continua} permite-nos a generalizar a obtenção da
densidade de probabilidade de transformações não monótonas que podem ser decompostas em
partes monótonas. Um exemplo clássico é a função $g(x) = x^2$ definida em $\mathbb{R}$.
Essa função não é monótona, porém suas partes definidas separadamente em $(-\infty, 0)$ e
$[0, \infty)$ serão.
A densidade da variável definida pela transformação poderá, portanto, ser obtida pelo
resultado a seguir

\begin{thm:fdp transformacao nao monotona continua}
Seja $X$ uma variável aleatória com f.d.p. igual a $f_X$.
Seja $Y=g(X)$ uma variável obtida pela transformação $g$ e com suporte $\Omega_Y$.
Se $\Omega_X$ admite partição $\{A_1, \ldots, A_n\}$ sendo possível definir a família
de funções $\{g_i: A_i \to \mathbb{R}\}_{i \in [n]}$ tal que
\begin{enumerate}
	\item[i] $g_i(x) = g(x)$
	\item[ii] $g_i(x)$ é monótona em $A_i$
	\item[iii] $\Omega_Y$ é o mesmo para todo $i \in [n]$.
	\item[iv] $g_i^{-1}$ tem derivada contínua em $\Omega_Y$ para todo $i \in [n]$.
\end{enumerate}
e $f_X$ contínua em cada $A_i$, então a f.d.p. de $Y$ é dada por
\[
	f_{Y}(y) = \begin{cases}
		\sum_{i=1}^nf_{X}(g_i^{-1}(y)) \cdot \left| \frac{d}{dy} g_i^{-1}(y)\right|, & \text{ se } y \in \Omega_{Y} \\
		0, \text{ caso contrário.}
	\end{cases}
\]
\end{thm:fdp transformacao nao monotona continua}
